\documentclass[11pt]{article}
\usepackage[letterpaper,margin=1in]{geometry}
\usepackage[T1]{fontenc}
\usepackage{lmodern,microtype}
\usepackage{amsmath,amssymb,amsthm,mathtools}
\usepackage{booktabs,array}
\usepackage[colorlinks=true,linkcolor=blue,citecolor=blue,urlcolor=blue]{hyperref}
\newtheorem{theorem}{Theorem}[section]
\newtheorem{lemma}[theorem]{Lemma}
\newtheorem{proposition}[theorem]{Proposition}
\newtheorem{corollary}[theorem]{Corollary}
\theoremstyle{definition}
\newtheorem{definition}[theorem]{Definition}
\newtheorem{assumption}[theorem]{Assumption}
\theoremstyle{remark}
\newtheorem{remark}[theorem]{Remark}
\newcommand{\R}{\mathbb R}
\newcommand{\Z}{\mathbb Z}
\newcommand{\N}{\mathbb N}
\newcommand{\G}{\mathcal G}
\newcommand{\norm}[1]{\left\lVert #1\right\rVert}
\newcommand{\abs}[1]{\left|#1\right|}
\newcommand{\ip}[2]{\langle #1,#2\rangle}
\DeclareMathOperator{\rank}{rank}
\DeclareMathOperator{\vol}{vol}
\DeclareMathOperator{\sign}{sign}
\DeclareMathOperator{\supp}{supp}
\newcommand{\poly}{\operatorname{poly}}
\title{Euclidean Steinitz Bounds for Graver Bases\\and Integer Programming\\[0.4em]
\large Conditional structural bounds and algorithmic consequences}
\author{AI-generated working draft\\Prepared for Robert Hildebrand}
\date{October 7, 2026}

\begin{document}
\maketitle
\begin{abstract}
Let $A\in\Z^{m\times n}$ have maximum Euclidean column norm $D$.
Assuming an absolute-constant Euclidean Steinitz bound $C\sqrt m$, we
prove $g_1(A)\le [K(1+D)]^m$, where $g_1$ is the maximum $\ell_1$-norm
of a Graver element and $K$ is absolute. The proof counts lattice points
in a Euclidean ball rather than its enclosing box. We obtain a
bounded-variable LP--IP proximity bound of $(m+1)[K(1+D)]^m$, give explicit
finite-state algorithms, and prove a two-level block-matrix bound
$G[K(1+D_A G)]^r$ with $G=[K(1+D_B)]^s$, up to integer rounding.
For columns with at most $s$ nonzeros of magnitude at most $\Delta$, the
general bound specializes to $[O(1+\sqrt{s}\Delta)]^m$.
A chain of equations supplies a matching $D^{\Omega(m)}$ lower bound.
We separately formulate a colorful consequence and explain why it does
not by itself improve the dependence on the number of blocks in four-block
integer programming. All main upper bounds are implications of an explicit
geometric hypothesis, as claimed in a September 2026 preprint. We do not
verify that preprint's analytic proof or claim a new best general-IP runtime.
\end{abstract}

\paragraph{Draft status.}
This is an exploratory AI-generated manuscript. The arguments below are
included in full where asserted as results. The Euclidean rearrangement
input is isolated as an assumption; the colorful section has a separate
input. Literature comparisons identify concrete earlier bounds, not an
exhaustive priority determination. The four-block application remains a
research direction, not a theorem of this paper. No computational solver
speedup or formal verification is claimed.

\section{Introduction}

Steinitz rearrangements give a useful way to study integer kernels: expand
a kernel vector into a zero-sum sequence of signed columns, reorder that
sequence into a bounded region, and use a repeated partial sum to extract
a smaller kernel vector. The geometry of the region matters as much as
the rearrangement bound. In particular, a Euclidean ball of radius
$O(\sqrt m D)$ contains at most $[O(1+D)]^m$ integer points. Enclosing
that ball in an axis-parallel box would introduce an avoidable factor
of order $m^{m/2}$.

Our purpose is to develop this observation systematically for integer
programming. Throughout, the number of rows is the dimension in which
Steinitz is applied; it is not the number of variables. The principal
matrix parameter is
\[
 D_2(A):=\max_{j\in[n]}\norm{A_{\cdot j}}_2.
\]
The familiar entry bound $\Delta=\max_{i,j}|A_{ij}|$ only implies
$D_2(A)\le\sqrt m\Delta$, and can conceal useful column geometry.

\paragraph{Results.}
Under Assumption~\ref{ass:ES}, we prove a Graver bound
$g_1(A)\le Q_m(D_2(A))$, where $Q_m(D)=[O(1+D)]^m$ is made explicit
below. For an optimal vertex of a linear relaxation with integral variable
bounds, we prove proximity at most $(m+1)Q_m(D)-1$. The proof handles
fractional coordinates by separating the rearranged walk into runs with
fixed lattice cosets. We then give a bounded-variable dynamic program, a
separable-convex improvement oracle, and a standard-form state graph.
For generalized $N$-fold matrices we apply the argument first to local
kernel vectors and then to their global images. The resulting structural
bound is independent of $N$ and the brick widths. Finally, we give a
balanced-subcollection consequence of a colorful Euclidean hypothesis.

\paragraph{Position relative to earlier work.}
Eisenbrand and Weismantel~\cite{EW} established the Steinitz approach to
proximity and integer-programming algorithms. Eisenbrand, Hunkenschr\"oder,
and Klein~\cite{EHK} used it to obtain the classical bound
$g_1(A)\le(2m\Delta+1)^m$ and a two-level block bound. Our arguments
retain their kernel-decomposition mechanism and change the geometry of
the counting step. Precise comparisons appear in Section~\ref{sec:comparison}.
Improvement over these particular bounds should not be confused with
improvement over all later algorithms.

\section{Geometric input and lattice counting}

\begin{assumption}[Euclidean rearrangement hypothesis]\label{ass:ES}
There is a fixed constant $C\ge1$ such that every finite sequence
$v_1,\ldots,v_L\in\R^m$ with $\sum_i v_i=0$ and
$\norm{v_i}_2\le D$ has a permutation $\pi$ satisfying
\[
 \norm{\sum_{i=1}^k v_{\pi(i)}}_2\le C\sqrt m D
 \qquad(0\le k\le L).
\]
This holds for every $m\ge1$ and $D\ge0$, with repetitions permitted.
\end{assumption}

The release entitled \emph{The Euclidean Steinitz--Bergstr\"om theorem}
\cite{ES2026} claims this hypothesis and a stronger prescribed-order
signing statement. We use the rearrangement assertion as a black box.
The algorithmic statements below are uniform once a fixed integer upper
bound on $C$ is supplied. They do not require an algorithm for finding
the asserted permutation.

Let $B_2^m$ denote the closed Euclidean unit ball and set
\[
 \kappa=\sqrt{2\pi e}\,(C+1),\qquad
 Q_m(D)=\left\lceil[\kappa(1+D)]^m\right\rceil.
 \tag{1}\label{eq:Q}
\]
Ceilings in this definition are only convenient explicit upper bounds;
implementations may replace these expressions by larger integer bounds.

\begin{lemma}[Uniform translated-ball count]\label{lem:count}
For every $a\in\R^m$ and $R\ge0$,
\[
 \#\big((a+\Z^m)\cap RB_2^m\big)
 \le \frac{\pi^{m/2}}{\Gamma(m/2+1)}(R+\sqrt m/2)^m
 \le \left[\sqrt{2\pi e}\left(\frac R{\sqrt m}+\frac12\right)\right]^m.
 \tag{2}\label{eq:count}
\]
In particular, the count is at most $Q_m(D)$ for $R=C\sqrt m D$.
\end{lemma}
\begin{proof}
Place a half-open unit cube centered at each of the counted lattice points.
The cubes have disjoint interiors, each has volume one, and their union
lies in $(R+\sqrt m/2)B_2^m$. This proves the first inequality.
The estimate $\vol(B_2^m)\le(2\pi e/m)^{m/2}$ proves the second.
Translation of the lattice has no effect on this argument.
\end{proof}

\begin{remark}[Rank and changes of coordinates]
The parameter $m$ can be reduced by deleting dependent rows; this preserves
the integer kernel and cannot increase $D_2$. For systems with right-hand
sides one must first check consistency of the deleted equations. An
arbitrary orthogonal change of coordinates does not preserve $\Z^m$.
Thus applying Lemma~\ref{lem:count} after such a rotation requires keeping
track of the rotated lattice. We do not silently replace it by $\Z^m$.
\end{remark}

\section{Graver elements and a lower bound}

For $u,v\in\Z^n$, write $u\sqsubseteq v$ if $u_i v_i\ge0$ and
$|u_i|\le|v_i|$ for every $i$. The Graver basis $\G(A)$ consists of the
$\sqsubseteq$-minimal nonzero members of $\ker(A)\cap\Z^n$.
Set $g_1(A)=\max\{\norm g_1:g\in\G(A)\}$, with value zero if the
integer kernel is trivial. Every nonzero integer kernel vector is a sum
of Graver elements conformal to it: repeatedly subtract a nonzero
$\sqsubseteq$-minimal kernel vector until the $\ell_1$-norm is zero.

\begin{theorem}[Euclidean-column Graver bound]\label{thm:graver}
Under Assumption~\ref{ass:ES}, every $A\in\Z^{m\times n}$ satisfies
\[
 g_1(A)\le Q_m(D_2(A)).
 \tag{3}\label{eq:graver}
\]
\end{theorem}
\begin{proof}
Fix $g\in\G(A)$ and let $L=\norm g_1$. Make a sequence containing
$|g_j|$ copies of $\sign(g_j)A_{\cdot j}$ for each nonzero coordinate.
The sum is zero and each vector has norm at most $D=D_2(A)$.
Use Assumption~\ref{ass:ES} and let $p_k$ be the reordered partial sums.
Then $p_0,\ldots,p_{L-1}$ are integer points of
$C\sqrt m D B_2^m$. They are distinct: if $p_i=p_j$ with
$0\le i<j\le L-1$, the intervening columns specify a nonzero integer
vector $h\sqsubseteq g$ with $Ah=0$ and
$\norm h_1=j-i<L$. There is no coordinate cancellation because all copies
of a given column have the sign prescribed by $g$. This contradicts
Graver minimality. Lemma~\ref{lem:count} now gives $L\le Q_m(D)$.
\end{proof}

\begin{corollary}[Entry and sparsity parameters]\label{cor:sparse}
If $|A_{ij}|\le\Delta$ and each column has at most $s\le m$ nonzeros,
then, under Assumption~\ref{ass:ES},
\[
 g_1(A)\le [O(1+\sqrt s\Delta)]^m.
\]
In particular the dense specialization is $[O(1+\sqrt m\Delta)]^m$.
\end{corollary}

\begin{proposition}[Exponential dependence cannot be removed]\label{prop:lower}
For every integer $q\ge2$ and $m\ge1$ there is an integer matrix
$A\in\Z^{m\times(m+1)}$ with at most two nonzeros per column such that
\[
 D_2(A)\le\sqrt{q^2+1},\qquad
 g_1(A)=1+q+\cdots+q^m\ge q^m.
\]
\end{proposition}
\begin{proof}
Use the equations $q x_i-x_{i+1}=0$, $1\le i\le m$.
Their integer kernel consists exactly of integer multiples of
$g=(1,q,\ldots,q^m)$. Hence $\G(A)=\{g,-g\}$.
The stated column bounds follow directly from the bidiagonal matrix.
\end{proof}

For this family, writing $D_0=\sqrt{q^2+1}$, the lower bound is at least
$(2D_0/\sqrt5)^m$. Thus the coefficient exponent $m$ in
Theorem~\ref{thm:graver} cannot be uniformly reduced for large $D$.
This does not assert matching constants or extremality for every $D$.

\subsection{An unconditional \texorpdfstring{$\ell_1$}{l1} comparison}
It is important to compare against classical Steinitz in norms other than
$\ell_\infty$. Let $D_1=\max_j\norm{A_{\cdot j}}_1$.
The arbitrary-norm Steinitz theorem (as stated in~\cite{EHK}) gives
partial sums in $mD_1B_1^m$.
Adding unit cubes enlarges the $\ell_1$ radius by at most $m/2$.
Since $\vol(B_1^m)=2^m/m!$, the same proof gives the unconditional bound
\[
 g_1(A)\le\frac{[2m(D_1+1/2)]^m}{m!}
 \le[2e(D_1+1/2)]^m.
 \tag{4}\label{eq:l1}
\]
One can round the right-hand side upward if desired. In particular,
$[O(1+s\Delta)]^m$ already follows for $s$-sparse columns. The
conditional Euclidean improvement replaces $s$ by $\sqrt s$ in this
comparison; single-exponential dependence for fixed $s,\Delta$ is not
a new qualitative conclusion. Matrices such as incidence matrices can
have much smaller Graver elements than either generic bound predicts.

\section{Proximity with integral variable bounds}\label{sec:proximity}

Consider
\begin{align*}
 \max\{c^\top x:Ax=b,\ \ell\le x\le u,\ x\in\R^n\}&\quad\text{(LP)},\\
 \max\{c^\top z:Az=b,\ \ell\le z\le u,\ z\in\Z^n\}&\quad\text{(IP)},
\end{align*}
where $A,b$ are integral and the finite entries of $\ell,u$ are integral.
We assume an optimal LP vertex $x^*$ and an optimal integer solution exist.
The objective vector may be rational. Bounds may be infinite when these
existence assumptions still hold.

\begin{theorem}[Bounded-variable proximity]\label{thm:proximity}
Under Assumption~\ref{ass:ES}, there is an optimal integer solution $z^*$
such that
\[
 \norm{z^*-x^*}_1\le(m+1)Q_m(D_2(A))-1.
 \tag{5}\label{eq:prox}
\]
More precisely, if $x^*$ has $f$ fractional coordinates, $m+1$ may be
replaced by $f+1$.
\end{theorem}
\begin{proof}
At a vertex, at most $\rank(A)\le m$ coordinates are strictly between
their bounds. Indeed, if there were more, their columns would admit a
nonzero dependence supported on those coordinates, giving a sufficiently
small feasible perturbation in both directions. A fractional coordinate
cannot equal an integral finite bound, so $f\le m$.

Choose an optimal integer solution $z$ minimizing $\norm{z-x^*}_1$.
Such a minimizer exists since a bounded ball contains only finitely many
integer points. Put $d=z-x^*$ and write
\[
 |d_j|=a_j+\theta_j,\qquad a_j=\lfloor|d_j|\rfloor,\quad 0\le\theta_j<1.
\]
Form a sequence of $a_j$ copies of $\sign(d_j)A_{\cdot j}$, together
with the vector $\theta_j\sign(d_j)A_{\cdot j}$ when $\theta_j>0$.
The sum is zero. Call the latter vectors exceptional; there are exactly
$f$ of them. Every vector has Euclidean norm at most $D=D_2(A)$.

Reorder the sequence by Assumption~\ref{ass:ES}. Exceptional vectors
split the integral vectors into at most $f+1$ consecutive runs, some
possibly empty. During any one run all partial sums lie in a fixed coset
of $\Z^m$ and in $C\sqrt m D B_2^m$. If a run contained $Q_m(D)$ or more
integral vectors, it would have at least $Q_m(D)+1$ partial sums,
including its initial point. Lemma~\ref{lem:count} would give two equal
sums. The intervening integral vectors specify a nonzero integer $h$
with $Ah=0$, the same coordinate signs as $d$, and $|h_j|\le a_j\le|d_j|$.

Both $x^*+h$ and $z-h$ lie coordinatewise between $x^*$ and $z$ and
are feasible for LP and IP, respectively. Optimality implies
$c^\top h\le0$ from LP and $c^\top h\ge0$ from IP. Thus $z-h$ is also
integer optimal, whereas
$\norm{z-h-x^*}_1=\norm{z-x^*}_1-\norm h_1$, a contradiction.

Every run therefore has at most $Q_m(D)-1$ integral vectors. Consequently
\[
 \norm{z-x^*}_1
 =\sum_j a_j+\sum_j\theta_j
 \le(f+1)(Q_m(D)-1)+f=(f+1)Q_m(D)-1.
\]
The weak inequality suffices, including the case $f=0$.
\end{proof}

\begin{remark}[Why fractional exceptions are kept]
Replacing the fractional residual by an arbitrary short integral
decomposition would require an additional argument and might increase
the column norm. The proof above never does this. It counts each
integral run in its own translated lattice, using the uniformity of
Lemma~\ref{lem:count}.
\end{remark}

The integrality gap is at most
$\norm c_\infty((m+1)Q_m(D)-1)$, by H\"older's inequality.
For rational data this statement applies after the usual integral
encoding of equality constraints; the resulting value of $D$ must be
computed for the actual encoded matrix.

\paragraph{Comparison with the Eisenbrand--Weismantel bound.}
For an entry bound $\Delta$, their bounded-variable proximity theorem
gives $m(2m\Delta+1)^m$~\cite{EW}. Since $D_2(A)\le\sqrt m\Delta$,
Theorem~\ref{thm:proximity} gives instead
\[
 (m+1)[O(1+\sqrt m\Delta)]^m.
\]
For $\Delta\ge1$, the displayed dimension factor is reduced from
$m^m$ to $m^{m/2}$, up to absolute constants to the power $m$ and
the leading polynomial factor. The stronger column-norm formulation
retains any additional saving from sparsity. This comparison is to that
specific theorem; later determinant-based proximity results must also
be considered, as discussed in Section~\ref{sec:comparison}.

\section{Explicit algorithmic consequences}\label{sec:algorithms}

The bounds are useful without constructing a Steinitz order. A dynamic
program may search a region guaranteed to contain an appropriate witness.
Nevertheless, a structural bound is not automatically a running-time
bound. We give explicit algorithms to make this distinction precise.
All counts below concern arithmetic operations on numbers of polynomial
encoding length, with the corresponding bit factors understood. Any
fixed constant used to instantiate Assumption~\ref{ass:ES} is part of
the algorithm specification, not something inferred from the input.

\subsection{A bounded-variable exact dynamic program}

\begin{proposition}\label{prop:boundedDP}
Suppose all variable bounds are finite integers. Given an optimal LP
vertex $x^*$, set $Q=Q_m(D)$ and $R=(m+1)Q+m$.
Under Assumption~\ref{ass:ES}, IP can be solved or declared infeasible
using
\[
 O\!\left(n(R+1)^2(2\lfloor DR\rfloor+1)^m\right)
 \tag{6}\label{eq:DPtime}
\]
arithmetic operations in addition to solving LP.
For nonzero integer $A$, the bit-time bound is
\[
 [O(1+D)]^{O(m^2)}\poly(n,L),
\]
where $L$ is the input length. LP infeasibility is handled first.
\end{proposition}
\begin{proof}
Let $w=\lfloor x^*\rfloor$ coordinatewise. Since the bounds are integral,
$w$ satisfies them, and $\norm{w-x^*}_1\le m$. If IP is feasible, finite
bounds imply existence of an optimum, and Theorem~\ref{thm:proximity}
gives an optimal solution $w+y$ with $\norm y_1\le R$.

Process coordinates in order. Store the best value of
$\sum_{j\le i}c_jy_j$ for every state
\[
 (i,t,k),\qquad t=\sum_{j\le i}A_{\cdot j}y_j\in\Z^m,
 \quad k=\sum_{j\le i}|y_j|\le R.
\]
Each transition chooses an integer $y_i\in[-R,R]$ satisfying
$\ell_i-w_i\le y_i\le u_i-w_i$. All reachable $t$ satisfy
$\norm t_2\le DR$, hence each coordinate has absolute value at most
$\lfloor DR\rfloor$. Initialize at $(0,0,0)$ and accept states with
$t=b-Aw$ and $k\le R$. This gives the stated state and transition count.
Every accepted solution is feasible, and an optimum is represented if
one exists. If $A=0$, independent coordinate optimization suffices.
\end{proof}

This intentionally simple algorithm has a quadratic-in-$m$ exponent.
It must not be advertised as a $[O(1+D)]^{O(m)}$ bounded-variable
algorithm. Faster augmentation or convolution routines would require
separate analysis.

\subsection{Separable convex objectives: an improvement oracle}

Let $F(x)=\sum_i f_i(x_i)$, where each $f_i$ is discretely convex on
its integer interval: $f_i(t+1)-f_i(t)$ is nondecreasing.
Suppose evaluations and comparisons are available exactly.

\begin{proposition}\label{prop:convex}
Given a feasible integer point $x$ with finite integer bounds, one can
find an improving feasible point or certify optimality for $F$ by searching
only directions $y$ with $Ay=0$ and $\norm y_1\le G$, where
$G=Q_m(D)$ under Assumption~\ref{ass:ES}.
The search uses
\[
 O\!\left(n(G+1)^2(2\lfloor DG\rfloor+1)^m\right)
\]
arithmetic operations and evaluations of objective differences.
\end{proposition}
\begin{proof}
Use the DP from Proposition~\ref{prop:boundedDP}, centered at $x$, with
budget $G$, target zero, and transition cost
$f_i(x_i+y_i)-f_i(x_i)$. Minimize this cost.
If $x$ is not optimal, let $z$ be better and decompose $z-x$ into Graver
elements $g^1,\ldots,g^t$ conformal to $z-x$. Each $x+g^j$ satisfies
the bounds. Discrete convexity, applied to monotone increments in each
coordinate, gives
\[
 F(z)-F(x)\ge\sum_{j=1}^t\big(F(x+g^j)-F(x)\big).
\]
Since the left-hand side is negative, at least one summand is negative.
Theorem~\ref{thm:graver} ensures that direction is searched.
\end{proof}

This is an exact optimality and improvement oracle, not a bound on the
number of iterations of unit-step augmentation. Repeatedly taking an
arbitrary improving step can have poor iteration complexity. In particular,
we do not infer a new general MINLP algorithm: separability and discrete
convexity are essential in this proposition.

\subsection{Standard form and a Euclidean tube}

Consider $Ax=b$, $x\in\Z^n_{\ge0}$, and put $B=\norm b_1$.
Assume $D>0$; for an integer matrix this implies $D\ge1$.
\begin{lemma}\label{lem:tube}
Under Assumption~\ref{ass:ES}, every feasible $x$ admits an ordering of
its multiset of columns whose partial sums lie in
\[
 [0,b]+C\sqrt m D B_2^m.
 \tag{7}\label{eq:tube}
\]
\end{lemma}
\begin{proof}
Append $M$ copies of $-b/M$ to the multiset, with
$M=\max\{1,\lceil\norm b_2/D\rceil\}$. All appended vectors have norm
at most $D$, and the total is zero. Apply Assumption~\ref{ass:ES} and
delete the appended vectors. A partial sum of the retained columns has
the form $p+(k/M)b$, where $\norm p_2\le C\sqrt m D$ and $0\le k\le M$.
\end{proof}

\begin{proposition}[Explicit tube graph]\label{prop:tubegraph}
Under Assumption~\ref{ass:ES}, standard-form feasibility has a finite
graph algorithm with at most
\[
 S\le (B+2)[K_0(1+D)]^m
 \tag{8}\label{eq:S}
\]
vertices and at most $nS$ directed edges, for an absolute constant $K_0$.
After construction, breadth-first search takes $O(nS)$ operations.
For a linear objective $c^\top x$, maximum-weight walks give either an
optimum or an unboundedness certificate in $O(nS^2)$ operations.
The graph can be constructed in $(B+2)[K_1(1+D)]^m\poly(n,m,L)$ time
for another absolute constant $K_1$.
\end{proposition}
\begin{proof}
For $B>0$, take the $B+1$ centers $a_k=(k/B)b$, $0\le k\le B$.
For $B=0$ take only the center zero. Adjacent centers have Euclidean
distance at most one. Let the vertices be all integer points in the
union of the balls of radius $T=C\sqrt m D+1$ about these centers.
This union contains the tube in \eqref{eq:tube}.
Lemma~\ref{lem:count}, with
$T/\sqrt m+1/2\le CD+3/2$, proves \eqref{eq:S} for a suitable $K_0$.
Include an edge $p\longrightarrow p+A_{\cdot j}$, labeled $j$, whenever
both endpoints occur. A walk from zero to $b$ gives a feasible nonnegative
integer solution by counting its labels. Lemma~\ref{lem:tube} proves
the converse. Edges of zero displacement are allowed.

For completeness, the vertex enumeration need not visit an enclosing
box. Enumerate coordinates of an integer point in a ball recursively,
retaining a prefix only if its squared distance in the assigned
coordinates is at most $T^2$. Every retained prefix has an integer
completion in the larger ball of radius $T+\sqrt m/2$: round the remaining
center coordinates to nearest integers. At any fixed depth distinct
prefixes have distinct such completions. Applying the cube argument
again bounds the total number of prefixes over all depths by
$(m+1)[\sqrt{2\pi e}(CD+2)]^m$. This proves the claimed construction
bound, including deduplication and edge lookup up to polynomial factors.
Radii may be enlarged to rational, polynomial-bit descriptions using a
fixed upper bound on $C$, without changing these estimates.

Give an edge labeled $j$ weight $c_j$. Every feasible solution yields a
walk of the same objective value, and every walk yields a feasible solution.
Consequently the graph represents exactly the attainable objective values.
Restrict to vertices reachable from zero and able to reach $b$. A positive
cycle there gives unbounded objective values; absent such a cycle, a
maximum-weight walk can be taken simple after deleting nonpositive
cycles. Bellman--Ford with negated weights gives the stated bound and
detects the relevant positive cycles.
\end{proof}

The dependence on $B$ is numerical, not logarithmic. The proposition is
an explicit geometric version of a state-graph argument, not a competitive
replacement for algorithms that avoid this dependence. It also does not
enforce upper bounds on label multiplicities, which is why bounded-variable
IP was treated separately.

\section{Two-level block matrices}\label{sec:block}

Let $r,s\ge1$ and consider
\[
 H=\begin{pmatrix}
 A_1&A_2&\cdots&A_N\\
 B_1&0&\cdots&0\\
 0&B_2&\cdots&0\\
 \vdots&\vdots&\ddots&\vdots\\
 0&0&\cdots&B_N
 \end{pmatrix},
 \qquad A_i\in\Z^{r\times t_i},\quad B_i\in\Z^{s\times t_i}.
\]
The brick widths $t_i$ need not be equal. Set
\[
 D_A=\max_{i,j}\norm{(A_i)_{\cdot j}}_2,
 \qquad D_B=\max_{i,j}\norm{(B_i)_{\cdot j}}_2.
\]

\begin{theorem}[Two-level Euclidean bound]\label{thm:block}
Under Assumption~\ref{ass:ES}, with $G=Q_s(D_B)$,
\[
 g_1(H)\le G\,Q_r(D_A G).
 \tag{9}\label{eq:block}
\]
More generally, $G$ can be any integer upper bound on
$\max_i g_1(B_i)$ with $G\ge1$, using Assumption~\ref{ass:ES} only in
the $r$-dimensional global step.
\end{theorem}
\begin{proof}
Take $g\in\G(H)$ and partition it into bricks $g_i$. Since $B_i g_i=0$,
write each nonzero brick as a conformal sum
$g_i=\sum_{j=1}^{q_i}h_{ij}$ with $h_{ij}\in\G(B_i)$.
The bound on $G$ gives $\norm{h_{ij}}_1\le G$.
The global images $v_{ij}=A_i h_{ij}$ satisfy
\[
 \sum_{i,j}v_{ij}=0,\qquad
 \norm{v_{ij}}_2\le\sum_k |(h_{ij})_k|\norm{(A_i)_{\cdot k}}_2
 \le D_A G.
\]
There can be no nonempty proper zero-sum subcollection of these images.
Such a subcollection would sum the corresponding local pieces to a
nonzero integer kernel vector $h\sqsubseteq g$ different from $g$.
The selected and omitted pieces are nonzero and conformal, so neither
can cancel internally. This would contradict Graver minimality.
The partial-sum argument of Theorem~\ref{thm:graver}, now in dimension
$r$, gives $\sum_i q_i\le Q_r(D_A G)$. Hence
$\norm g_1=\sum_{i,j}\norm{h_{ij}}_1\le GQ_r(D_A G)$.
\end{proof}

The formula does not depend on the number of bricks. It also permits
sharper application-specific local bounds in place of $Q_s(D_B)$.
For example, if a local block has a small explicitly known Graver basis,
there is no reason to discard that information.

\paragraph{Comparison with the classical two-level calculation.}
For an entry bound $\Delta\ge1$, the corresponding bound in
\cite[Lemma 3]{EHK} is
\[
 G_{\rm old}(2r\Delta G_{\rm old}+1)^r,
 \qquad G_{\rm old}=(2s\Delta+1)^s.
 \tag{10}\label{eq:blockold}
\]
Our dense substitution $D_B\le\sqrt s\Delta$, $D_A\le\sqrt r\Delta$
gives, with absolute constants displayed schematically,
\[
 G\,[K(1+\sqrt r\Delta G)]^r,
 \qquad G\le[K(1+\sqrt s\Delta)]^s.
 \tag{11}\label{eq:blockdense}
\]
Ignoring only factors $K^{O(rs+r+s)}$, the dimension factors in
\eqref{eq:blockold} are $s^{s(r+1)}r^r$ whereas those in
\eqref{eq:blockdense} are $s^{s(r+1)/2}r^{r/2}$; the coefficient exponent
in both comparisons is $s(r+1)+r$. This is a comparison of two structural
bounds obtained by the same two-level decomposition, not an assertion
about the best available complete $N$-fold algorithm.

If local columns have at most $a$ nonzeros and global columns at most
$b$, replace $\sqrt s$ and $\sqrt r$ by $\sqrt a$ and $\sqrt b$ in
\eqref{eq:blockdense}. However, the global images $A_i h_{ij}$ need not
themselves be sparse. The proof uses their norm bound $D_A G$, not an
unjustified assertion that they retain the original support size.

\paragraph{Algorithmic meaning.}
Theorem~\ref{thm:block} bounds the total $\ell_1$-size of a Graver
augmentation independently of $N$. It can be used as a structural input
where an augmentation routine has an explicitly stated dependence on
this size. Running the generic DP of Section~\ref{sec:algorithms} on
all $r+Ns$ rows would lose the block advantage. A faster block-aware
implementation and its full iteration bound are not supplied here.

\section{Colorful sequences and what they do not yet prove}\label{sec:colorful}

The ordinary zero-sum hypothesis is sufficient for all preceding results.
This section uses the following distinct input, also claimed as a
consequence in the released preprint~\cite{ES2026}.

\begin{assumption}[Colorful Euclidean hypothesis]\label{ass:colorful}
For $k$ sequences $v_{j,1},\ldots,v_{j,L}\in\R^d$, $1\le j\le k$,
with $\norm{v_{j,i}}_2\le D$ and $\sum_{j,i}v_{j,i}=0$, there are
permutations $\pi_j$ such that
\[
 \norm{\sum_{j=1}^k\sum_{i=1}^{t}v_{j,\pi_j(i)}}_2
 \le C_{\rm col}\sqrt d D\qquad(0\le t\le L),
\]
where $C_{\rm col}$ is absolute and independent of $k,L$.
\end{assumption}

The source obtains $C_{\rm col}=2C$ from its stronger signing result
and matrix transference. We do not claim that
Assumption~\ref{ass:colorful} has been proved here from
Assumption~\ref{ass:ES} alone. Define
\[
 Q^{\rm col}_d(D)=
 \left\lceil[\sqrt{2\pi e}(C_{\rm col}+1)(1+D)]^d\right\rceil.
\]

\begin{definition}
An array of integer vectors $v_{j,i}$ with $k$ rows and $L$ entries per
row, whose total sum is zero, is \emph{balanced-indecomposable} if there
are no subsets $I_j\subseteq[L]$ of a common cardinality
$t\in\{1,\ldots,L-1\}$ such that
$\sum_j\sum_{i\in I_j}v_{j,i}=0$.
\end{definition}

\begin{proposition}\label{prop:balanced}
Under Assumption~\ref{ass:colorful}, a balanced-indecomposable array
with $\norm{v_{j,i}}_2\le D$ has
\[
 L\le Q^{\rm col}_d(D).
 \tag{12}\label{eq:balanced}
\]
\end{proposition}
\begin{proof}
Let $p_t$ be its synchronous partial sums after the row permutations.
The $L$ integer points $p_0,\ldots,p_{L-1}$ lie in
$C_{\rm col}\sqrt d D B_2^d$. A repetition $p_a=p_b$, $a<b<L$,
selects $b-a$ entries from each row with total zero, contrary to the
definition. The translated-ball count proves the assertion.
\end{proof}

The bound controls the number of rounds $L$, not the total number $kL$
of vectors. This distinction is essential. With $L=1$, any number of rows
can have total zero, and the balanced-indecomposability condition has no
proper positive number of rounds to exclude. Thus the proposition cannot
alone remove the number of colors from a Graver norm bound.

\begin{proposition}[Centering when the total is nonzero]\label{prop:center}
Under Assumption~\ref{ass:colorful}, an arbitrary array of the same
shape and vector bound admits row permutations for which, writing
$T=\sum_{j,i}v_{j,i}$,
\[
 \norm{\sum_{j=1}^k\sum_{i=1}^t v_{j,\pi_j(i)}-\frac tL T}_2
 \le 2C_{\rm col}\sqrt d D\qquad(0\le t\le L).
\]
\end{proposition}
\begin{proof}
Subtract the mean $\mu_j=L^{-1}\sum_i v_{j,i}$ from every entry in
row $j$. Each centered vector has norm at most $2D$ and the total is
zero. Apply Assumption~\ref{ass:colorful} and then restore the means.
\end{proof}

\paragraph{Four-block IP remains an application to complete.}
Oertel, Paat, and Weismantel~\cite{OPW} develop colorful rearrangements
for block-structured proximity. B\'ar\'any~\cite{Barany} later improves
the arbitrary-norm row-permutation bound to $4d-2$.
Therefore a comparison only against the earlier $40d^5$ bound would
overstate what is new. With entrywise magnitude $M$, Euclidean conversion
gives $D\le\sqrt d M$, so the new rearrangement radius is $O(dM)$,
the same coordinate scale as the later arbitrary-norm bound. An additional
saving can come from ball counting, provided the proof really supplies
integer partial sums or a controlled number of lattice cosets.

To obtain a full four-block theorem, one must verify three further facts
inside that decomposition: the Euclidean norms of the derived vectors,
the precise balanced selection that yields an admissible kernel move,
and the number of exceptional shifts or rounds that need separate counting.
Neither Proposition~\ref{prop:balanced} nor
Proposition~\ref{prop:center} establishes these facts. We therefore make
no claim here of an improved polynomial exponent in the number of blocks.

\section{Comparisons and limits of the conclusions}\label{sec:comparison}

\paragraph{Dense general-IP algorithms.}
Jansen and Rohwedder~\cite{JR} already give algorithms with the
$\sqrt m\Delta$ base, and formulate their analysis using hereditary
discrepancy. Thus the dense specialization of our bound does not by
itself improve the known general-IP runtime. Our tube algorithm also
retains a numerical dependence on $\norm b_1$. Its role is to show
constructively how the geometric hypothesis can certify a smaller state
space, rather than to establish a best running time.

\paragraph{Determinant-based proximity.}
For full row rank, let $\Delta_m(A)$ be the largest absolute determinant
of an $m\times m$ submatrix. In nonnegative standard form,
Lee, Paat, Stallknecht, and Xu~\cite{LPSX} prove proximity less than
\[
 3m^2\log(2\sqrt m\,\Delta_m(A)^{1/m})\,\Delta_m(A).
\]
Hadamard gives $\Delta_m(A)\le D_2(A)^m$. Their result can therefore
be stronger than our direct bound on the overlapping formulation.
Theorem~\ref{thm:proximity} includes integral upper and lower bounds,
but we do not infer novelty merely from this difference without checking
the applicable extensions. It is a complete implication of the geometric
hypothesis and a convenient direct proof, not a universal improvement.

\paragraph{What the sparse comparison establishes.}
For fixed $s,\Delta$, both \eqref{eq:l1} and Corollary~\ref{cor:sparse}
give $2^{O(m)}$. For growing $s$, the displayed Euclidean upper bound
replaces a base proportional to $s\Delta$ by one proportional to
$\sqrt s\Delta$. The constants are absolute and inherited from the
geometric hypothesis. This comparison does not establish that every
matrix with those parameters has Graver elements near either upper bound.

\paragraph{Practical interpretation.}
Smaller structural bounds can reduce a certified search region, but hidden
constants, state representation, preprocessing, and instance structure
all matter. No implementation or benchmark in this manuscript supports
a claim of faster commercial branch-and-bound or cutting-plane solvers.
Likewise, the convex improvement oracle does not extend to arbitrary
nonconvex objectives. A useful experimental follow-up would compare
actual reachable states in Euclidean and coordinate-box neighborhoods
for selected bounded-column instances.

\section{Further questions}

The most immediate theoretical task is a comprehensive novelty audit of
the Euclidean-column Graver bound, including bounds expressed through
determinants, circuits, discrepancy, and other column norms. The lower
example proves the necessity of exponential growth in $m$ and the
coefficient scale, but does not settle the best absolute base.

A second task is a full block-aware augmentation analysis using
Theorem~\ref{thm:block}. The structural bound should be substituted into
an explicitly chosen algorithm, including step lengths, feasibility
initialization, and the number of augmentation phases. Doing so would
determine whether the dimension-factor improvement survives all other
terms in the runtime.

A third task is the four-block program described in
Section~\ref{sec:colorful}. Its main difficulty is preserving the correct
integer or translated-lattice structure through the colorful decomposition.
An improvement in a rearrangement radius alone does not settle it.

Finally, one can ask for anisotropic versions: choose a norm adapted to
the columns and count the actual image lattice in the resulting body.
Any such result must control both the body and the lattice after
preconditioning. Optimizing only the column norm while forgetting the
lattice density would give an invalid improvement.

\appendix
\section{Dependency and verification ledger}

The logical dependencies of the stated results are as follows.
\begin{center}
\begin{tabular}{>{\raggedright\arraybackslash}p{0.40\textwidth}>{\raggedright\arraybackslash}p{0.48\textwidth}}
\toprule
Result & External geometric input\\
\midrule
Translated-ball count; chain lower bound & None\\
Unconditional $\ell_1$ comparison & Classical arbitrary-norm Steinitz theorem\\
Graver bound; proximity; state algorithms & Assumption~\ref{ass:ES}\\
Two-level block bound & Assumption~\ref{ass:ES} locally and globally, or a supplied local bound\\
Balanced-array bound; centering & Assumption~\ref{ass:colorful}\\
Improved four-block IP theorem & Not proved in this draft\\
\bottomrule
\end{tabular}
\end{center}

The proofs were developed as implications of the stated assumptions.
The analytic proof in the September 2026 release has not been independently
verified here. The release is a mutable online source; the bibliography
records its dated directory and the access date. Before submission, an
archived revision or commit should be fixed, author metadata verified,
and any available corrections checked. No Lean verification is included.
Finite examples or numerical experiments, if added, cannot replace the
universal arguments above. The repository accompanying this draft records
compilation and any small-instance checks separately.

\begin{thebibliography}{9}
\raggedright

\bibitem{ES2026}
OpenAI math preprint release,
\emph{The Euclidean Steinitz--Bergstr\"om theorem},
dated directory September 24, 2026.
\href{https://github.com/openai/math/tree/main/preprints/The-Euclidean-Steinitz-Bergstrom-theorem-September-24-2026}{Repository source}.
The statements used here are in \texttt{build/introduction.tex} and
\texttt{build/consequences.tex}. Accessed October 7, 2026.
Repository attribution is used here without assigning individual authorship.

\bibitem{EW}
F. Eisenbrand and R. Weismantel,
\emph{Proximity results and faster algorithms for Integer Programming
using the Steinitz Lemma},
\href{https://arxiv.org/abs/1707.00481}{arXiv:1707.00481}, 2017;
version 3, 2019.

\bibitem{EHK}
F. Eisenbrand, C. Hunkenschr\"oder, and K.-M. Klein,
\emph{Faster Algorithms for Integer Programs with Block Structure},
\href{https://arxiv.org/abs/1802.06289}{arXiv:1802.06289}, 2018.

\bibitem{JR}
K. Jansen and L. Rohwedder,
\emph{On Integer Programming, Discrepancy, and Convolution},
\href{https://arxiv.org/abs/1803.04744v4}{arXiv:1803.04744v4}, 2022.

\bibitem{LPSX}
J. Lee, J. Paat, I. Stallknecht, and L. Xu,
\emph{Improving proximity bounds using sparsity},
\href{https://arxiv.org/abs/2001.04659}{arXiv:2001.04659}, 2020.

\bibitem{OPW}
T. Oertel, J. Paat, and R. Weismantel,
\emph{A Colorful Steinitz Lemma with Application to Block-Structured
Integer Programs},
\href{https://arxiv.org/abs/2201.05874v2}{arXiv:2201.05874v2}, 2022.

\bibitem{Barany}
I. B\'ar\'any,
\emph{A matrix version of the Steinitz lemma},
\href{https://arxiv.org/abs/2308.10102v2}{arXiv:2308.10102v2}, 2024.

\end{thebibliography}
\end{document}
